\begin{center}
	\large{\textbf{OPTIMASI PORTOFOLIO INVESTASI BERBASIS EKONOFISIKA DAN GAME THEORY MENGGUNAKAN \textit{VARIATIONAL QUANTUM EIGENSOLVER}}}
\end{center}
\begin{minipage}{5cm}
\begin{align*}
&\text{Nama} \qquad &&: \text{PUTRA NAUFAL TABASSAM} \\
&\text{NRP} \qquad &&: \text{5001221120} \\
&\text{Departemen} \qquad &&: \text{FISIKA} \\
&\text{Pembimbing} \qquad &&: \text{Dr. Heru Sukamto, M.Si.} \\
\end{align*}
\end{minipage} \\
\textbf{Abstrak} 
\par Model optimasi \textit{Mean-Variance} sering digunakan dalam investasi pasa modal secara kuantitatif. Namun model ini memiliki keterbatasan dalam mengatasi kompleksitas pasar, seperti korelasi antar aset yang ekstrem dan volatilitas yang tinggi, terutama ketika terjadi peristiwa ekonomi penting. Untuk mengatasi tantangan tersebut, tugas akhir ini mengusulkan metode hibrida dengan mengintegrasikan prinsip ekonofisika, teori permainan, dan komputasi kuantum variasional sehingga dapat meningkatkan akurasi serta efisiensi pengambilan keputusan investasi pasar modal. Permasalahan alokasi aset ditransformasi menjadi formulasi \textit{Quadratic Unconstrained Binary Optimization} (QUBO) yang dipetakan ke dalam sistem Hamiltonian Ising. Kerangka \textit{Exact Potential Game} (EPG) diterapkan untuk memodifikasi parameter medan lokal ($h_i$) berdasarkan utilitas marginal strategis, sementara parameter kopling ($J_{ij}$) menangkap dinamika interaksi risiko antar aset. Optimasi dilakukan menggunakan algoritma \textit{Variational Quantum Eigensolver} (VQE) dengan inisialisasi \textit{warm-start} berbasis \textit{Nash Equilibrium} dari kerangka EPG yang sama untuk memitigasi fenomena \textit{Barren Plateaus} dan mempercepat konvergensi pada sirkuit kuantum era NISQ. Hasil simulasi menggunakan data saham indeks LQ45 periode 2021--2023 menunjukkan bahwa strategi GT-VQE mampu menghasilkan performa yang lebih baik dengan nilai \textit{Sharpe Ratio} mencapai 1,0107 pada sistem jumlah aset $N=4$. Selain itu, integrasi teori permainan juga terbukti mengefisiensikan daya komputasi lewat pemilihan kedalaman sirkuit yang lebih dangkal ($L \leq 3$) dan dapat menjaga stabilitas pertumbuhan modal di tengah fluktuasi pasar walaupun terjadi peristiwa ekonomi ekstrem.

\begin{minipage}{1cm}
	\begin{align*}
	\text{\textbf{Kata kunci}}:\hspace{2mm}&\text{\textbf{\textit{Ekonofisika, Teori Permainan, Model Ising,}}} \\
	&\text{\textbf{Optimasi Portofolio, VQE}}
	\end{align*}
\end{minipage}

\newpage
\thispagestyle{empty}
\vspace*{\fill}
    \begin{center}
	Halaman ini sengaja dikosongkan
    \end{center}
\vspace*{\fill}
\pagebreak
